\section{Real Data Application}\label{SC:RESULT}\label{sec:results}
The experiments apply the SAD components defined in Section
\ref{SC:METHOD} to two insurance-claims benchmarks. The evaluation compares
the retained architectures, individual classifiers, a frozen-model control
and two joint threshold settings.

\subsection*{Carclaims dataset}
Carclaims contains 15,420 records, 32 original predictors and 923 recorded
frauds (5.99\%) \citep{dataset}. The original descriptive frequency table
is retained and regenerated from the supplied data in Table \ref{tab:frequency}.
\begin{table}[tbp]\centering\footnotesize
\caption{Frequency table for selected categorical variables in Carclaims. Original descriptive variables are retained; percentages are recomputed from the supplied file.}\label{tab:frequency}
\begin{tabular}{llrr}\toprule Feature & Category & Percent & Frequency\\\midrule
Sex & Female & 15.69 & 2,420\\
 & Male & 84.31 & 13,000\\
MaritalStatus & Divorced & 0.49 & 76\\
 & Married & 68.90 & 10,625\\
 & Single & 30.38 & 4,684\\
 & Widow & 0.23 & 35\\
VehicleCategory & Sedan & 62.72 & 9,671\\
 & Sport & 34.75 & 5,358\\
 & Utility & 2.54 & 391\\
DriverRating & 1 & 25.58 & 3,944\\
 & 2 & 24.65 & 3,801\\
 & 3 & 25.19 & 3,884\\
 & 4 & 24.58 & 3,791\\
FraudFound & No & 94.01 & 14,497\\
 & Yes & 5.99 & 923\\
\bottomrule\end{tabular}\end{table}


\subsection{Additional dataset and replication scope}\label{sec:seconddata}
Insurance Claims is available through the Kaggle dataset published by
\texttt{mexwell} \citep{insuranceclaims}, the source identified for the
second file supplied in the research archive. The analyzed file contains
1,000 records, 38 original predictors and 247 recorded frauds (24.70\%).
Its schema, record count and contents differ from Carclaims.
The differently named \texttt{fraud\_oracle.csv} is a duplicate
representation of Carclaims and is excluded from the additional benchmark.

Because the feature spaces differ, preprocessing and models are fitted
separately for each dataset under the same SAD protocol. This experiment
assesses replication of the workflow across benchmarks. The Kaggle source
documents access to the second dataset; the original collecting insurer
and its full real-versus-synthetic provenance remain unverified.

\begin{table}[tbp]\centering
\caption{Input datasets before predictor screening.}\label{tab:data}
\begin{tabular}{lrrr}\toprule
Dataset & Records & Predictors & Recorded fraud\\\midrule
Carclaims & 15,420 & 32 & 923 (5.99\%)\\
Insurance Claims & 1,000 & 38 & 247 (24.70\%)\\
\bottomrule\end{tabular}\end{table}

\subsection{Algorithm specifications}\label{sec:specifications}
The number of partitions is $m=4$, with proportions
$p_1=0.50$, $p_2=0.30$, $p_3=0.10$ and $p_4=0.10$.
Simple-A uses $N=9$ initial detectors; Simple-B uses $N=10$ by including
EAD as a vote; Multilayer retains $N=9$ and applies EAD to $Y_1$.
All three retain $M=8$ supervised learners:
LR, SVM, XGB, RF, NB, KNN, AdaBoost and ANN.
Seeds 1234--1243 and the four-part schedule are shared across variants.
The main operating point is $(T_{AD},T_{ML})=(3,4)$; the stricter point $(5,6)$
is a complete additional Multilayer training trajectory.
These are fixed illustrative settings, not a preregistered optimum.
The data had been examined during development; the second benchmark is not
described as prospectively untouched external data.
SAD hyperparameters are unchanged between datasets.
Only the PSO comparator follows the predefined inner training-only search.
Random partitions are stratified using the outcomes for split construction
only. Within each trajectory, label access follows Algorithms 2 and 3.

\paragraph{Reconciliation with the original exports.}
All reported performance values are recomputed from the corrected
implementation. The revision uses training-only preprocessing and SMOTE, corrected
entropy/index alignment, and outcome access only after review selection.
Non-alert samples are verified rather than assumed genuine. The previous
exports did not consistently designate fraud as the positive class.
Consequently, their reported recall/F1 and high overall accuracy are not
reused as evidence of fraud detection performance. These corrections are
necessary to evaluate the retained SAD architectures consistently.

\FloatBarrier
\subsection{Architecture performance and workload}
\begin{table}[!htbp]\centering\footnotesize
\caption{Fraud-class performance and total review fraction; mean $\pm$ sample SD across ten seeds. All metrics pool $D_2$--$D_4$ within a seed. AP denotes average precision.}\label{tab:architectures}
\setlength{\tabcolsep}{2pt}
\begin{tabular}{lrrrrr}\toprule
Architecture & Precision & Recall & F1 & AP & Review\\\midrule
\multicolumn{6}{l}{\textit{Carclaims}}\\
Simple-A & $0.144 \pm 0.009$ & $0.480 \pm 0.039$ & $0.221 \pm 0.014$ & $0.151 \pm 0.009$ & $0.280 \pm 0.012$\\
Simple-B & $0.146 \pm 0.006$ & $0.489 \pm 0.044$ & $0.224 \pm 0.009$ & $0.153 \pm 0.008$ & $0.281 \pm 0.016$\\
Multilayer & $0.144 \pm 0.008$ & $0.407 \pm 0.047$ & $0.212 \pm 0.014$ & $0.140 \pm 0.007$ & $0.253 \pm 0.016$\\
Multilayer-static & $0.141 \pm 0.008$ & $0.427 \pm 0.046$ & $0.211 \pm 0.013$ & $0.135 \pm 0.009$ & $0.264 \pm 0.017$\\
\multicolumn{6}{l}{\textit{Insurance Claims}}\\
Simple-A & $0.491 \pm 0.015$ & $0.586 \pm 0.078$ & $0.533 \pm 0.038$ & $0.487 \pm 0.029$ & $0.368 \pm 0.033$\\
Simple-B & $0.510 \pm 0.038$ & $0.624 \pm 0.038$ & $0.560 \pm 0.028$ & $0.491 \pm 0.039$ & $0.377 \pm 0.028$\\
Multilayer & $0.465 \pm 0.035$ & $0.544 \pm 0.101$ & $0.497 \pm 0.043$ & $0.452 \pm 0.048$ & $0.365 \pm 0.050$\\
Multilayer-static & $0.435 \pm 0.042$ & $0.546 \pm 0.104$ & $0.479 \pm 0.041$ & $0.442 \pm 0.034$ & $0.386 \pm 0.067$\\
\bottomrule\end{tabular}
\end{table}
\begin{table}[!htbp]\centering\footnotesize
\caption{Complementary metrics under the same main protocol. High accuracy alone is insufficient for these imbalanced datasets.}\label{tab:architectures_aux}
\setlength{\tabcolsep}{2pt}
\begin{tabular}{lrrrr}\toprule
Architecture & Specificity & Balanced acc. & ROC-AUC & Accuracy\\\midrule
\multicolumn{5}{l}{\textit{Carclaims}}\\
Simple-A & $0.818 \pm 0.013$ & $0.649 \pm 0.017$ & $0.777 \pm 0.010$ & $0.797 \pm 0.011$\\
Simple-B & $0.818 \pm 0.016$ & $0.653 \pm 0.016$ & $0.777 \pm 0.008$ & $0.798 \pm 0.013$\\
Multilayer & $0.846 \pm 0.016$ & $0.626 \pm 0.018$ & $0.752 \pm 0.017$ & $0.819 \pm 0.013$\\
Multilayer-static & $0.834 \pm 0.018$ & $0.631 \pm 0.017$ & $0.744 \pm 0.018$ & $0.810 \pm 0.015$\\
\multicolumn{5}{l}{\textit{Insurance Claims}}\\
Simple-A & $0.800 \pm 0.025$ & $0.693 \pm 0.028$ & $0.738 \pm 0.015$ & $0.747 \pm 0.009$\\
Simple-B & $0.800 \pm 0.035$ & $0.712 \pm 0.020$ & $0.742 \pm 0.019$ & $0.756 \pm 0.024$\\
Multilayer & $0.792 \pm 0.045$ & $0.668 \pm 0.035$ & $0.707 \pm 0.026$ & $0.730 \pm 0.020$\\
Multilayer-static & $0.761 \pm 0.070$ & $0.654 \pm 0.032$ & $0.694 \pm 0.024$ & $0.708 \pm 0.035$\\
\bottomrule\end{tabular}
\end{table}

On Carclaims, Simple-B has the highest mean F1 among the four architecture variants (0.224). Simple-B and updating Multilayer attain AP 0.153 and 0.140, respectively; Multilayer's mean F1 is 0.212. The corresponding review fractions are 0.281 for Simple-B and 0.253 for Multilayer.

On Insurance Claims, Simple-B has the highest mean F1 among the four architecture variants (0.560). Simple-B and updating Multilayer attain AP 0.491 and 0.452, respectively; Multilayer's mean F1 is 0.497. The corresponding review fractions are 0.377 for Simple-B and 0.365 for Multilayer.

These differences concern both predictive performance and the records acquired for training. The architectures are therefore compared as complete review policies, without holding the review budget constant.

Figures \ref{fig:performance}--\ref{fig:workload} summarize pooled performance,
partition-wise detection and verification workload.
\paperfigure{02_architecture_performance}{Fraud recall, F1 and AP for four
architectures. Points and bars show mean and sample SD across ten random splits,
pooling $D_2$--$D_4$ within each seed.}{fig:performance}{1}
\paperfigure{03_partition_progress}{Partition-wise F1 and recall.
Later partitions contain different records; these are not learning curves
on a fixed test set.}{fig:partitions}{1}
\paperfigure{04_workload}{Partition-wise alert rate and total review rate,
including the 10\% audit of non-alerted claims. Lines and points show means
over ten random splits. The panels distinguish alerted claims from the
full set submitted for verification.}{fig:workload}{1}

\subsection{Classifier comparisons and ranking curves}\label{sec:comparators}
The eight component classifiers are evaluated individually on the same
verified training records acquired by Simple-B. Two additional comparators
use those same records: balanced random forest (200 trees with within-tree
balancing and no preceding SMOTE) and PSO-XGBoost adapted from the insurance
fraud study of \citet{Ding2025}. The latter is a recent fraud-motivated
comparator; it is not an exact reproduction of the published system.

PSO selects hyperparameters once using only verified $D_1$ records and
three-fold stratified cross-validation scored by AP. SMOTE is fitted inside
each training fold. Four particles and three iterations search 60--160 trees,
depth 2--6, learning rates $10^{-2}$--$10^{-0.7}$ and row/column subsampling
0.6--1. Inertia is 0.7 and both attraction coefficients are 1.4.
Selected values are fixed in subsequent rounds. The comparison tests
classifier choice under a common acquisition policy, not ten independent
end-to-end acquisition systems.
\begin{table}[!htbp]\centering\footnotesize
\caption{Conditional classifier comparison using Simple-B's verified records. The two columns per benchmark report fraud F1 and AP. PSO-XGB is the explicitly adapted comparator.}\label{tab:baselines}
\setlength{\tabcolsep}{2pt}
\begin{tabular}{lrrrr}\toprule
Method & Carclaims F1 & Carclaims AP & Insurance F1 & Insurance AP\\\midrule
LR & $0.221 \pm 0.007$ & $0.143 \pm 0.008$ & $0.508 \pm 0.021$ & $0.461 \pm 0.035$\\
SVM & $0.217 \pm 0.008$ & $0.137 \pm 0.007$ & $0.501 \pm 0.037$ & $0.455 \pm 0.036$\\
XGB & $0.076 \pm 0.017$ & $0.159 \pm 0.009$ & $0.551 \pm 0.033$ & $0.497 \pm 0.036$\\
NB & $0.155 \pm 0.022$ & $0.101 \pm 0.016$ & $0.500 \pm 0.022$ & $0.460 \pm 0.022$\\
RF & $0.039 \pm 0.015$ & $0.161 \pm 0.006$ & $0.537 \pm 0.032$ & $0.511 \pm 0.046$\\
KNN & $0.172 \pm 0.008$ & $0.101 \pm 0.005$ & $0.413 \pm 0.013$ & $0.296 \pm 0.026$\\
AdaBoost & $0.216 \pm 0.012$ & $0.144 \pm 0.006$ & $0.558 \pm 0.030$ & $0.512 \pm 0.036$\\
ANN & $0.188 \pm 0.011$ & $0.140 \pm 0.010$ & $0.404 \pm 0.033$ & $0.360 \pm 0.046$\\
BalancedRF & $0.231 \pm 0.012$ & $0.168 \pm 0.007$ & $0.595 \pm 0.025$ & $0.497 \pm 0.029$\\
PSO-XGB-adapted & $0.116 \pm 0.073$ & $0.158 \pm 0.010$ & $0.599 \pm 0.026$ & $0.522 \pm 0.032$\\
Simple-B & $0.224 \pm 0.009$ & $0.153 \pm 0.008$ & $0.560 \pm 0.028$ & $0.491 \pm 0.039$\\
\bottomrule\end{tabular}
\end{table}

For Carclaims, BalancedRF leads this comparator set in both mean F1 (0.231) and AP (0.168). For Insurance Claims, PSO-XGB-adapted leads this comparator set in both mean F1 (0.599) and AP (0.522). These rankings are conditional on Simple-B's acquisition policy and the stated training settings. Differences in tuning budgets limit comparisons with fully optimized implementations of the published methods.

Figures \ref{fig:baselines} and \ref{fig:curves} show the classifier comparison
and the architectures' ranking performance on the predefined split.
\paperfigure{05_baseline_comparison}{Same-acquisition classifier comparison.
All comparators use the Simple-B verified-record policy.
PSO-XGB is an adapted search; BalancedRF balances within trees.
Error bars show sample SD.}{fig:baselines}{1}
\paperfigure{06_roc_pr}{ROC and precision--recall curves for predefined seed
1234, pooling $D_2$--$D_4$. These are representative split curves,
not averages over independent populations. Dashed lines are no-skill
references.}{fig:curves}{1}

\subsection{Ablation study and threshold sensitivity}\label{sec:ablation}
Comparing Simple-A with Simple-B examines adding a flat entropy vote;
comparing Simple-B with Multilayer examines replacing flat entropy voting
with candidate-local entropy gating. These comparisons also change which
records are reviewed and therefore are not matched-budget causal estimates.
For the sequential ablation, Multilayer-static freezes the exact $D_1$
models and predicts the same $D_2$--$D_4$ records without retraining.
Its initialization, preprocessing and initial training set are identical
to those of updating Multilayer.
\begin{table}[!htbp]\centering\footnotesize
\caption{Paired seed-wise change from static to updating Multilayer: mean $\pm$ sample SD. These are descriptive paired differences, not confidence intervals.}\label{tab:ablation}
\setlength{\tabcolsep}{2pt}
\begin{tabular}{lrrr}\toprule
Dataset & $\Delta$ F1 & $\Delta$ AP & $\Delta$ review\\\midrule
Carclaims & $0.0005 \pm 0.0081$ & $0.0049 \pm 0.0052$ & $-0.0109 \pm 0.0070$\\
Insurance Claims & $0.0182 \pm 0.0223$ & $0.0098 \pm 0.0211$ & $-0.0216 \pm 0.0304$\\
\bottomrule\end{tabular}
\end{table}

For Carclaims, sequential updating changes mean F1 by +0.0005, AP by +0.0049 and review fraction by -0.0109 relative to the frozen model. For Insurance Claims, sequential updating changes mean F1 by +0.0182, AP by +0.0098 and review fraction by -0.0216 relative to the frozen model. The Carclaims F1 change is small relative to its split-to-split variation. These descriptive results do not establish statistical superiority. Because updating affects later acquisition as well as prediction, the ablation measures the complete update policy. The two threshold settings are evaluated through separate complete trajectories.

\paperfigure{07_ablation_deltas}{Paired seed-wise changes, in percentage
points, from frozen to updating Multilayer. Positive F1/AP changes favor
updating; review-rate differences measure changed workload. Black diamonds
and bars show mean and sample SD. Panels use separate vertical
ranges.}{fig:ablation}{1}
\begin{table}[!htbp]\centering\footnotesize
\caption{Multilayer sensitivity at two joint threshold settings; mean $\pm$ sample SD.}\label{tab:thresholds}
\setlength{\tabcolsep}{2pt}
\begin{tabular}{lrrrr}\toprule
Dataset / thresholds & Recall & F1 & AP & Review\\\midrule
Carclaims / (3,4) & $0.407 \pm 0.047$ & $0.212 \pm 0.014$ & $0.140 \pm 0.007$ & $0.253 \pm 0.016$\\
Carclaims / (5,6) & $0.111 \pm 0.029$ & $0.123 \pm 0.021$ & $0.127 \pm 0.007$ & $0.142 \pm 0.010$\\
Insurance Claims / (3,4) & $0.544 \pm 0.101$ & $0.497 \pm 0.043$ & $0.452 \pm 0.048$ & $0.365 \pm 0.050$\\
Insurance Claims / (5,6) & $0.300 \pm 0.120$ & $0.369 \pm 0.118$ & $0.440 \pm 0.059$ & $0.231 \pm 0.051$\\
\bottomrule\end{tabular}
\end{table}

\paperfigure{08_threshold_sensitivity}{Main $(3,4)$ and strict $(5,6)$
Multilayer points from separately rerun trajectories.
The comparison does not isolate the two thresholds independently.}{fig:thresholds}{1}

\subsection{Computation and reproducibility}
Predictions retain original row ID, split, seed, architecture, score, votes,
alert, review mask and recorded outcome. Manifests record software versions
and code/data hashes. Numerical verification recomputes metrics from these
records; figures and LaTeX tables use the same results.
The executed implementation is Python, using scikit-learn and the associated
packages listed in the environment record \citep{Pedregosa2011}.
Original R scripts and exports are archived; revised results come from Python.
A single R/\texttt{ggplot2} script renders both benchmarks from the saved
results. The empirical figures use the same visual conventions and share
axis limits within each comparison. Performance axes span 0--80\% and
workload axes 0--55\%, encompassing the displayed values and error bars;
ROC and precision--recall curves retain the full 0--1 range. Numerical
tables report fractions on the 0--1 scale. Error bars denote sample
standard deviations across splits, not confidence intervals.

\begin{table}[!htbp]\centering\footnotesize
\caption{Wall time in seconds, mean $\pm$ sample SD. Architecture rows sum supervised fitting at $D_1,D_2,D_3$; the shared anomaly stage is reported separately and excludes preprocessing.}\label{tab:runtime}
\setlength{\tabcolsep}{2pt}
\begin{tabular}{lrr}\toprule
Operation & Carclaims & Insurance Claims\\\midrule
Simple-A & $15.94 \pm 0.79$ & $1.69 \pm 0.10$\\
Simple-B & $17.67 \pm 1.21$ & $1.83 \pm 0.09$\\
Multilayer & $10.84 \pm 0.79$ & $1.30 \pm 0.06$\\
Shared anomaly stage & $2.38 \pm 0.23$ & $0.39 \pm 0.03$\\
\bottomrule\end{tabular}
\end{table}

\paperfigure{09_runtime}{Supervised fitting time across three training
rounds, excluding shared initial anomaly calculation, exports and comparator
tuning. Times describe this hardware and implementation; the panels use
different time ranges.}{fig:runtime}{1}


\paragraph{Data and code availability.}
The source datasets are identified in \citet{dataset} and
\citet{insuranceclaims}. The analyzed files, data dictionaries, executable
analysis and plotting scripts, prediction-level outputs and instructions
for reproducing the tables and figures are publicly available in the
project GitHub repository.\footnote{\url{https://github.com/Mehrab-Atighi/Anomaly-Detection-Algorithm-for-Enhanced-Fraud-Identification-in-the-Insurance-Sector}}
The archived analysis uses the supplied dataset versions; source hashes
are included to identify these files unambiguously.
